%%%
%%%
%%%
\section{Overview}
%%%
\subsection{Main principles}\label{subsec_main_principles}
The following, general development principles have been retained by the founding developer team.
%%%
\paragraph*{Minimality.}
For each \texttt{sw<Model>} class, a single .h file and several .cpp files if needed.
%%%
\paragraph{Subsidiarity.}
Classes, functions, methods are defined at the lowest possible level. For instanec, the time loop (for a transient model), initialization of parameters and variables, writing of results are model-dependent. Consequently, they are defined individually for each \texttt{sw<Model>} class. 


%%%
\subsection{Glossary}\label{subsec_glossary}
The following definitions are used hereafter.
\begin{description}
\item [Model] (in the context of \swSolver). A model is a simplified representation of reality. It means:
\begin{itemize}
\item a set of state variables (generally, but not necessarily, conserved variables)
\item a set of governing equations describing the behaviours of these variables
\item the definitions of the fluxes, source terms, etc. representing the various processes at work
\item a set of space- and/or time-dependent parameters
\item a number of numerical techniques managing the solution of the governing equations (e.g. Godunov, Muscl-Evr, etc.), tailored specifically to the model under consideration.
\end{itemize}
\item [Process] A process is attached to a specific part of the governing equations. So far, three types are recognized in the \swSolver nomenclature:
\begin{itemize}
\item Hyp for the hyperbolic part of the equations
\item Parab for the parabolic part (diffusion)
\item SourceTerms for the source terms part (RHS of the governing equations)
\end{itemize}
\item [Type] A particular type of source term. Examples: Wind, Friction, PrecipitationInfiltration, SingularHeadLoss
\item [sub-type] Within a given type: a specific formulation for this type. Example: for Type = Friction, sub-type = Manning.
\item [Method] A specific numerical method used for the time integration of the governing equations. Examples: Godunov, Muscl-Evr.
\end{description}
%%%
%
\section{Directories and files}\label{sec_sourceFileNameStructure}
See also \ref{sec_namingRules}.
\subsection{Overview of folder structure}\label{subsec_overview_folder_structure}

\begin{tabular}{| p{6cm} | p{8cm} |}
\hline
\texttt{src/} & the general folder for all source code files \\
\hline
\texttt{src/common} & source files used by several apps \\
\hline
\texttt{src/sw2dConverter} & source files that are specific to the \swConverter executable \\
\hline
\texttt{src/sw2dSolver} & source files that are specific to the \swSolver executable \\
\hline
\texttt{src/sw2dSolver/common} & general purpose source files  in the \swSolver executable \\
\hline
\texttt{src/sw2dSolver/models} & folder containing the various model subfolders \\
\hline
\texttt{src/sw2dSolver/models/common} & folder containing the files that are used by several models \\
\hline
\texttt{src/sw2dSolver/models/<model>} & folder containing the files for Model \texttt{<model>} (e.g. liswes, swes, sp). \\
\hline
\texttt{src/sw2dSolver/models/swParent.*} & header + source file for the \texttt{<Parent>} class from which all models inherit. \\
\hline
\texttt{src/sw2dSolver} & source files that are specific to the \swZone executable \\
\hline
\end{tabular}\\
%%
%%
\subsection{in \texttt{<Model>} folders}
\paragraph{General rule.}
A .cpp file name has the following structure:
\begin{lstlisting}
 sw<Model><Process><Type><Sub-type><Method>.cpp
 \end{lstlisting}
  All arguments between \texttt{< >} are optional, however the order of these arguments (even empty) should be preserved.
There is only 1 .h file for a given model : \texttt{sw<Model>.h}. All .cpp files falling within the context of this \texttt{<Model>} use the header file \texttt{sw<Model>.h} and are not allowed to use the header from another model.
%%%
\paragraph{Example\#1:} \texttt{swSwesHypGodunov.cpp} (uses \texttt{swSwes.h}): 
\begin{itemize}
\item \texttt{<Model> = Swes} (the Shallow water equations)
\item \texttt{<Process> = Hyp} (the hyperbolic part of the equations is being dealt with)
\item \texttt{<Type>, <Sub-type>} are empty (meaningless for hyperbolic terms)
\item \texttt{<Method> = Godunov} (the first-order Godunov scheme is used to solve the hyperbolic part of the shallow water equations)
\end{itemize}
%%%
%%%
\paragraph{Example\#2:} \texttt{swSwesSourceFrictionChezy} (uses \texttt{swSwes.h}):
\begin{itemize}
\item \texttt{<Model> = Swes} (the Shallow water equations)
\item \texttt{<Process> = Source} (the source term part of the equations is being dealt with)
\item \texttt{<Type> = Friction} within the source term, the specific process bein handled is friction
\item \texttt{<Sub-type> = Chezy} the friction term is computed using Chezy's approach.
\item \texttt{<Method>} is empty because there is only one method used for the time integration, so there is no need to define a specific keyword for it (it would be different if e.g. an explicit + implicit + RK2 methods were availble for instance).
\end{itemize}
%
%
%
%
\section{Class, method and variable naming rules}\label{sec_namingRules}
%
%
%
\subsection{Class naming}\label{subsec_classNamingRules}
SW2D classes always start with the \texttt{sw} prefix. This prefix is followed by a number of keywords, each starting with a capital letter for improved readability.

%
%
\paragraph{common/ classes.}

%
%
\paragraph{models/ classes.}

%
%
%
\subsection{Method naming}
Wihtin a given class, a given method name always starts with a lowercase character. When the method is made of several words, any new word starts with an uppercase character. Example:
\texttt{computeTimeStep}
%
%
%
\subsection{variable naming}
\begin{itemize}
\item Variable names are made of lowercase characters. One of the very few exceptions are L (and its mirror R) to denote the left- and right-hand side attributes of interfaces, beause most code editors do not allow for an easy distinction between lowercase L and 1 (unit number).
\item When a variable is made of several words, each word is separated from the preceding by an underscore (\_) character.
\item When the variable refers to a specific attribute for an object, the variable name should always follow a descending hierarchy through object properties. 

Example:
\texttt{edge\_midpoint\_coordinate\_x[i]} refers to the $x-$ coordinate of the midpoint of an edge. The logic is the following:
\begin{itemize}
\item an \text{edge} has a number of attributes, such as its length, normal unit vector, etc. Here, the particular attribute under consideration is its midpoint, hence \texttt{\_midpoint}.
\item one of the attributes of the midpoint is its set of coordinates, hence \texttt{\_coordinate}
\item the particular coordinate of interest is the $x-$ coordinate, hence \texttt{\_x}.
\end{itemize}
\end{itemize}

%
%
%
%
\section{Coding rules}

Each class should be described in a dedicated set of 2 files: \texttt{<class>.h} and \texttt{<class>.cpp}.
%
%
%
\subsection{Headers}

Header (.h) files have the following structure.
\begin{enumerate}
	\item \texttt{\# pragma once} instruction.
	\item \texttt{\# include< >} for system/language classes (e.g. \texttt{<vector>}, \texttt{<string>}), in alphabetical order.
	\item \texttt{\# include " "} instructions for sw2d classes(e.g. "swIO.h"), in alphabetical order .
	\item \texttt{class} instruction defining the local classes.
	\item For each class
	\begin{enumerate}
	\item public methods (alphabetical order),
	\item public attributes (alphabetical order),
	\item Protected methods (alphabetical order),
	\item Protected attributes (alphabetical order),
	\item Private methods (alphabetical order),
	\item Private variables (alphabetical order).
	\end{enumerate}
\end{enumerate}
%
%
%
\subsection{CPP files}

C++ (.cpp) have the following structure.

\begin{enumerate}
	\item \texttt{\# include} instructions for external (e.g. math, IO) libraries,in alphabetical order,
	\item \texttt{\# include} instructions for sw2d libraries, in alphabetical order,
	\item Constructor method for the class,
	\item Destructor method for the class,
	\item other methods, in alphabetical order, irrespective of their public/protected/private nature.
\end{enumerate}
%
%
%
\subsection{Code commenting}
%
%
\paragraph{Reserved keywords.}
The following keywords are reserved to help the developer team to track issues in the code. They are used as entries in the Doxygen-generated documentation. They may appear in .h and/or .cpp files.
\begin{itemize}
\item \texttt{/** @bug: <text> **/} Signalling a bug.
\item \texttt{/** @todo: <text> **/} Tasks that are incomplete/still to be done.
\item \texttt{/** @why: <text> **/} Request for information.

\end{itemize}

%
%
%
%
\subsection{Doxygen fields}
%
%
%
\paragraph{Header files.}
The attributes of the class are briefly documented in the header file next to their declaration. The following bracketing should be used for the purposes of Doxygen: \texttt{/**<  **/}

Example: 
\begin{verbatim}
     std::vector<double> z_bottom;   /**< [i]: bottom level of cell with Id i **/
\end{verbatim}
%
%
\paragraph{C++ files.}
For each method, the purpose, parameters and return type of the method should be described using Doxygen formatting:
\begin{itemize}
\item \texttt{/**} to start the Doxygen commenting sequence
\item \texttt{@brief} followed by the decription of the method (on the same line)
\item \texttt{@param} name of the variable, followed by its decription (1 line per parameter). Can be in or out.
\item \texttt{@return} meaning of the variable returned by the method (in the case where only one variable is expected)
\item \texttt{*/} to close the Doxygen commenting sequence
\end{itemize}
Example:
\begin{verbatim}
/**
 * @brief Computes fluxes at boundary edges for the SWES models
 * @param[in] celL: Id of cell on the left-hand side of the interface 
 * @param[in] celR: Id of cell on the right-hand side of the interface 
 * @param[in] qL  : normal unit discharge in the cell on the LHS of the interface 
 * @param[in] qR  : normal unit discharge in the cell on the RHS of the interface 
 * @param[in] rL  : tangent unit discharge in the cell on the LHS of the interface 
 * @param[in] rR  : tangent unit discharge in the cell on the RHS of the interface 
 * @param[in] side: the side of the interface on which the boundary is located ('L' or 'R')
 * @param[in] zL  : free surface elevation in the cell on the LHS of the interface
 * @param[in] zR  : free surface elevation in the cell on the RHS of the interface
 * @param[out] F1 : mass flux in the direction normal to the interface
 * @param[out] F2 : momentum flux in the direction normal to the interface
 * @param[out] F3 : momentum flux in the direction tangent to the interface 
 * @param[out] s2R: geometric source term on the RHS of the interface
 */

\end{verbatim}
